% Created 2017-11-09 Thu 16:45
% Intended LaTeX compiler: pdflatex
\documentclass[11pt]{article}
\usepackage[utf8]{inputenc}
\usepackage[T1]{fontenc}
\usepackage{graphicx}
\usepackage{grffile}
\usepackage{longtable}
\usepackage{wrapfig}
\usepackage{rotating}
\usepackage[normalem]{ulem}
\usepackage{amsmath}
\usepackage{textcomp}
\usepackage{amssymb}
\usepackage{capt-of}
\usepackage{hyperref}
\usepackage[margin=0.80in]{geometry}
\usepackage[backend=biber, style=ieee]{biblatex}
\usepackage{float}
\usepackage[super,negative]{nth}
\usepackage[capitalise]{cleveref}
\usepackage{pst-node,transparent,ragged2e}
\usepackage[inline]{enumitem}
\addbibresource{/home/nuk3/course/training/csir/novellasers/bibliography/bibliography.bib}
\DeclareFieldFormat[inproceedings]{citetitle}{\textit{#1}}
\DeclareFieldFormat[inproceedings]{title}{\textit{#1}}
\DeclareFieldFormat[misc]{citetitle}{#1}
\DeclareFieldFormat[misc]{title}{#1}
\renewcommand*{\bibpagespunct}{%
\ifentrytype{inproceedings}
{\addspace}
{\addcomma\space}}
\AtEveryCitekey{\ifuseauthor{}{\clearname{author}}}
\AtEveryBibitem{\ifuseauthor{}{\clearname{author}}}
\renewcommand\thesection{(\alph{section})}
\renewcommand\thesubsection{}
\PassOptionsToPackage{hyperref,x11names}{xcolor}
\hypersetup{colorlinks=true,citecolor=blue,filecolor=cyan,linkcolor=black,urlcolor=blue}
\SpecialCoor
\author{Nyameko Lisa}
\date{27 October 2017}
\title{IPSP044 - Essential Trade Mark Law\\\medskip
\large Assignment 03 - 788684}
\hypersetup{
 pdfauthor={Nyameko Lisa},
 pdftitle={IPSP044 - Essential Trade Mark Law},
 pdfkeywords={},
 pdfsubject={},
 pdfcreator={Emacs 27.0.50 (Org mode 9.1.2)},
 pdflang={English}}
\begin{document}

\maketitle
\addvspace{110pt}
\centering{
\pnode(0.5\textwidth,-0.5\textheight){thisCenter}
\rput(thisCenter){%\transparent{0.25}
\includegraphics[width=2.7in]{/home/nuk3/course/llb/wipo-unisa/UNISACoatofArms.eps}}}

\justifying
\addvspace{110pt}
\section*{Declaration}
\label{sec:org4ffabf4}
I know that plagiarism is to use someone else’s work and pass it off as my own.\\
I know that plagiarism is wrong.\\
I confirm that this assignment is my own work.\\
I have acknowledged in the bibliography accompanying the assignment all the sources that I have used.\\
I have not directly copied without acknowledgement anything from the Internet or from any other source.\\
I have indicated every quotation and citation in a footnote or bracket linked to that quotation.\\
I have not allowed anyone else to copy my work and to pass it off as their own work.\\
I understand that if any unacknowledged copying whatsoever appears in my assignment I will receive zero per cent for the assignment.\\
I am aware of the UNISA policy on plagiarism and understand that disciplinary proceedings can be instituted against me by UNISA if I contravene this policy.\\
I indicate my understanding and acceptance of this declaration by
entering my name hereunder:
\begin{itemize}
\item Name: \textbf{Nyameko Lisa} (Student Number: \textbf{7874-909-3})
\end{itemize}

\subsection*{NOTE}
\label{sec:org9feef80}
Please note that footnotes will be denoted as \footnote{This is a footnote.} and will
appear at the bottom of the page.\\
References will be denoted by \cite{rsa78_patents_act} and will appear at the end of the document.
\newpage

\section{Advise Aroma Company whether the mark qualifies for registration as a trade mark. [25]}
\label{sec:orge7e4e5f}
As per the definitions of [section 2(1)]\cite{rsa93_tm_act} and given that the
Republic is a Paris Convention Member State, [article 15(1)]\cite{wto17_trips},
[article 1(2)]\cite{wipo83_paris_conve_protect_ip}, it follows that \emph{`TROPIC
GLOW'} will be understood to be a \textbf{mark}, in that it is either a name or sign
capable of being represented graphically. Moreover it constitutes a \textbf{trade
mark}, as it has been proposed by Aroma Company, in relation to distinguishing
their trade or use from other goods or services connected in the course of trade
regarding \emph{`soaps, perfumery, essential oils and cosmetics'}, \cite{wipo57_ncl}. Therefore it
follows as per [section 9(1)]\cite{rsa93_tm_act}, that TROPIC GLOW qualifies for
registration as trade mark, subject to limitations, in relation to use within
those limitations, and thus the unregistrability conditions of the
provisions of [paragraphs (1), (4), (5), (8), (9), (10) or (11) of section 10]\cite{rsa93_tm_act},
cannot be applied.\\

As per [section 2(4)]\cite{rsa93_tm_act}, \emph{`use'} or \emph{`proposed use'} shall be
construed as the use of a visual representation of the mark [section
2(2)(a)]\cite{rsa93_tm_act}, as pertaining to physical form or other relation in
terms of goods[section 2(3)(a)]\cite{rsa93_tm_act}, or as pertaining to the use
thereof in any relation to the performance of any services\footnote{Also referred to
as service marks, [article 1(2)]\cite{wipo83_paris_conve_protect_ip}} [section
2(3)(b)]\cite{rsa93_tm_act}. Provided that no others claim priority to the TROPIC
GLOW mark at the  of application for registration, then as per [section
9(2)]\cite{rsa93_tm_act}, TROPIC GLOW shall be considered to be capable of
distinguishing within the meaning described above. Thus also negating the
unregistrability conditions of the provisions of [paragraphs (2), (3) or (6) of section 10]\cite{rsa93_tm_act}\\

As per [section 11(1)]\cite{rsa93_tm_act}, the TROPIC GLOW mark is proposed to be
registered in respect of goods falling in a particular class, namely class 3,
hence [section 10(4)]\cite{rsa93_tm_act} cannot be applied. Moreover TROPIC GLOW
does not include or consist of the name or representation of a person [section
12]\cite{rsa93_tm_act}, hence there would be no need to furnish the registrar with
additional legal consent from any parties. There would be no need for the
registrar to impose honest concurrent use conditions and limitations [section
14(1)]\cite{rsa93_tm_act}, as the mark TROPIC GLOW does not offend against the
provisions of [paragraphs (6), (14), (15) or (17) of section
10]\cite{rsa93_tm_act}.\\

It has already been demonstrated that TROPIC GLOW contains matter which is
indeed capable of distinguishing within the meaning of [section
9]\cite{rsa93_tm_act}, specifically in this sense as a composite mark even though
it's separate components ``TROPIC'' and ``GLOW'' may individually not be capable
of distinguishing. Hence registration may be subject to disclaimer or memorandum
[section 15]\cite{rsa93_tm_act}, and the registrar may require Aroma Company to
disclaim any right to the exclusive use of either ``TROPIC'' and / or ``GLOW''
seperately and apart from TROPIC GLOW, [section 15(a)]\cite{rsa93_tm_act}. One
such example exists in South African case, \citetitle{vivier2000_cadbury_v_beacon}
\cite{vivier2000_cadbury_v_beacon}, where registration of the Beacon \emph{Liquorice
Allsorts} mark was rectified and subject to the following disclaimers
\begin{enumerate}
\item ``Registration of the Beacon \textit{Liquorice Allsorts} trade mark shall give no right to the exclusive use
of the sweet device [the plate], separately and apart from the mark.''
\item ``Registration of this mark shall also give no right to the exclusive use
of the name \textit{Liquorice Allsorts}, separately and apart from the mark.''
\end{enumerate}

Provided that Aroma Company and its use of TROPIC GLOW does not inherently [section 10(12)]\cite{rsa93_tm_act}
or from the manner of its use [section 10(13)]\cite{rsa93_tm_act}, likely
deceive, cause confusion or offend any class of person, it cannot be deemed
unregistrable. Moreover provided that TROPIC GLOW does not constitute an
identical or similar trade mark registration [section 10(14)]\cite{rsa93_tm_act}
or prior application [paragraphs (15) and (16) of section 10]\cite{rsa93_tm_act},
where use of TROPIC GLOW would likely amount to deceit or confusion in relation
to the goods and services of a different proprietor or applicant, respectively,
without their consent; the registrar cannot refuse registration on those
grounds. Lastly provided TROPIC GLOW does not constitute an identical or similar
mark already registered and well known in the Republic, [section
10(17)]\cite{rsa93_tm_act}, the registration cannot be refused on the grounds of
unfair competition through the dilution of good will.

\section{Advise Aroma Company whether any of these prior applications / registrations will present an obstacle to the registration of its marks. [25]}
\label{sec:orgd8fc367}
\subsection{Registration Number 90/2390 TROPICAL in class 25}
\label{sec:org4e25a4d}
As per the provisions of [section 10(17)]\cite{rsa93_tm_act}, the proprietors of
the TROPICAL trade mark could argue that their registered trade mark is
well-known in the Republic [section 35(1)]\cite{rsa93_tm_act} and that
registration of the TROPICAL GLOW trade mark without their consent, albeit in a
different class to that of TROPICAL, would likely amount to unfair advantage or
be detrimental to the distinctive character or repute of the TROPICAL trade
mark.
\subsection{Registration Number 96/4782 TROPICANA in class 3}
\label{sec:org1493c1d}
As the provisions of [Article
$6^{bis}$]\cite{wipo83_paris_conve_protect_ip}, the registered trade
mark TROPICANA can be considered a \emph{`well-known'} mark, moreover it follows that
the proprietors of TROPICANA are either \begin{enumerate*}[label=(\roman*)]
\item nationals of the Republic [section 35(1)(a)]\cite{rsa93_tm_act}, or \item
domiciled in or have a real and effective industrial or commercial establishment
in the Republic [section 35(1)(b)]\cite{rsa93_tm_act}, which TROPICANA is
considered to be.\end{enumerate*} Given that the TROPICANA trade mark has been
in use extensively since 1996 in relation to lipsticks, class 3 in respect of
cosmetics, \cite{wipo57_ncl}, knowledge of the trade mark to the public in the relevant sector can
be assumed [section 35(1A)]\cite{rsa93_tm_act}.\\

If Aroma Company can demonstrate that they have used the TROPIC GLOW trade mark
in the Republic, at some time before 1996 when the TROPICANA trade mark was in
use, claiming priority to the TROPIC GLOW mark [Article
4]\cite{wipo83_paris_conve_protect_ip}, then Aroma Company shall be entitled to the
``Saving of vested rights'' provision [section 36(2)]\cite{rsa93_tm_act} and
regardless of whether or not TROPICANA is a well-known mark entitled to
protection under the Paris Convention, nothing in the Act shall permit the
proprietors of TROPICANA to interfere with or restrain use / registration of
TROPIC GLOW. Otherwise the proprietors of TROPICANA may argue that the TROPIC
GLOW trade mark is unregistrable on the grounds that essential parts thereof
constitute reproduction or imitation of their mark TROPICANA, which is both
well-known [section 35(1)]\cite{rsa93_tm_act} and entitled to protection under the
Paris Convention [section 35(3)]\cite{rsa93_tm_act}
\cite{wipo83_paris_conve_protect_ip}. Therefore TROPICANA would indeed present an
obstacle to the registration of TROPIC GLOW as they trade under cosmetics
\emph{(lipsticks)} \cite{wipo57_ncl}, which are indeed identical or similar goods and services, and the
proprietors of TROPICANA may argue that use of TROPIC GLOW is likely to cause
deception or confusion [section 10(6)]\cite{rsa93_tm_act}.\\

Unless Aroma Company can demonstrate that the TROPIC GLOW trade mark constitutes
\emph{``honest concurrent use''} [section 14(1)]\cite{rsa93_tm_act}, or unless they can
obtain consent from the proprietors of the TROPICANA trade mark, then Aroma
Company may face obstacles in that the proprietors may oppose registration on
the grounds that TROPIC GLOW would likely confuse or deceive [section
10(14)]\cite{rsa93_tm_act}, or constitute unfair advantage or competition through
the dilution of the good-will of TROPICANA [section 10(17)]\cite{rsa93_tm_act}.\\

Should Aroma Company proceed to use the TROPIC GLOW mark, without consent from
the proprietors of the TROPICANA mark, they may find themselves as respondents
to infringement proceedings brought against them by the proprietors of
TROPICANA, as the provisions of [section 34]\cite{rsa93_tm_act}, and as per the
instance in South African case law, \citetitle{corbett97_mcd_v_joburgers}
\cite{corbett97_mcd_v_joburgers}, where the respondent Joburgers infringed on the
plaintiff's trade mark \emph{McDonald's}, through the unlawful use of the similar
\emph{MacDonald's} mark.

\subsection{Application Number 2000/764 TROPIC LIGHT in class 3}
\label{sec:org9aeb511}
Without demonstration of prior use [section 9(2)]\cite{rsa93_tm_act} or consent
from the proprietors of the TROPIC LIGHT trade mark, they may oppose
registration of the TROPIC GLOW mark, on the grounds that it constitutes a mark
so similar to their prior application that it would be likely to cause deceit of
confusion [section 10(15)].\\

Moreover, provided that the proprietors of TROPIC LIGHT have advertised their
application for registration in the prescribed manner [section
17]\cite{rsa93_tm_act}, then Aroma Company may themselves file an opposition to
that registration [section 21]\cite{rsa93_tm_act}, should they be able to
demonstrate that registration of TROPIC LIGHT is contrary to the existing rights
of Aroma Company and its subsequent registration of the TROPIC GLOW mark
[section 10(16)]\cite{rsa93_tm_act}.\\

Lastly, in light of all of the above considerations, the registrar may permit
registration of the TROPIC GLOW trade mark but possibly subject to a number of
disclaimers pertaining to the TROPICAL, TROPICANA and TROPIC LIGHT trade marks
respectively, as the provisions of [section 15]\cite{rsa93_tm_act}.

\section{What must Aroma Company establish in order to succeed with any infringement action against Beaver Company? [25]}
\label{sec:org9e0b1ea}

As per the provisions of [section 34(1)(a)]\cite{rsa93_tm_act}, Aroma Company may
institute infringement proceedings against Beaver Company for its unauthorised
use of the TROPICLEAN mark in class 3 in respect of cosmetic products \emph{(hair
lotions)} \cite{wipo57_ncl}, and it's similarity to the registered trade mark TROPIC GLOW. As per
the provisions of [section 34(5)]\cite{rsa93_tm_act} before Aroma Company can
institute any proceedings in terms of infringement, they will be required to
give written notice of such intention to all users concerned in the register
and said users are also entitled to intervene and recover any damages they also
may have suffered as a result of the infringement. In order to succeed with any
infringement proceedings against Beaver Company, Aroma Company must demonstrate
(i) probability of deception of confusion, (ii) unauthorised use and (iii)
infringement in relation to goods or services for which trade mark is registered.

\subsection{Existence of a probability of deception or confusion}
\label{sec:org12d680d}
Aroma Company is required to demonstrate that the allegedly infringing mark
TROPICLEAN is indeed used as a trade mark for the purpose of indicating the
origin of Beaver Company's goods. This follows directly from the similarity
between the infringing trade mark of the hair products TROPICLEAN and that of
Aroma Company's TROPIC GLOW cosmetics, \cite{wipo57_ncl}. Moreover Beaver Company has commenced
marketing the infringing product in similar packing to that used by Aroma
Company and thus create the impression of a material link between the Beaver
Company's TROPICLEAN and Aroma Company's TROPIC GLOW, which is likely to cause
deception or confusion.

\subsection{Alleged infringing use must be unauthorised}
\label{sec:org4c8a3bf}
Aroma Company need only demonstrate that they have not provided any consent to
Beaver Company to use their trade mark, through either assignment, transmission, sale,
hypothecation or license [sections (39), (40) and (41)]\cite{rsa93_tm_act}. Specifically they must
establish that use of the infringing mark does not constitute permitted use, and
that Beaver Company are not considered to be registered users, as per the
provisions of [section 38]\cite{rsa93_tm_act}.

\subsection{Alleged infringing use must be in relation to goods or services for which the TROPIC GLOW is registered}
\label{sec:org0468540}
TROPIC GLOW is registered in class 3 in respect of cosmetics products.
TROPICLEAN is marketed as a hair products, which is indeed a cosmetic product.
This follows naturally.

\subsection{Alleged infringing use not protected by provisions of [section 34(2)]}
\label{sec:org7372c47}
Aroma Company must demonstrate that TROPICLEAN does not constitute bone fide use
of a (i) name of person or their predecessors, (ii) nor their place of business,
[section 34(2)(a)]\cite{rsa93_tm_act} - of which it would be easy to demonstrate
that TROPICLEAN is not; a description or indication of the nature or intended
purpose of said goods, [section 34(2)(b-c)]\cite{rsa93_tm_act} - again TROPICLEAN
is neither; no consent has been approved, [section 34(2)(d)]\cite{rsa93_tm_act};
the similar packaging of TROPICLEAN to that used by Aroma Company in TROPIC GLOW
is for deceptive and not utilitarian purposes [section
34(2)(e)]\cite{rsa93_tm_act}; and TROPICLEAN is not a registered trade mark,
[section 34(2)(g)], simple search of the trade marks register will resolve this.

\subsection{Relief that Aroma Company may pursue}
\label{sec:org392c869}
Aroma Company should be advised that it is not easy to succeed in any infringement
proceedings, as the example in South African case law between
\citetitle{harms07_verimark_v_bmw} \cite{harms07_verimark_v_bmw} demonstrates. Where
the Supreme Court of Appeal found in favour of the initial respondent Verimark
against the plaintiff that initiated the infringement proceedings Bayerische
Motoren Werke Aktiengesellschaft (BMW), where BMW failed on appeal in an attempt to seek
an interdict for unfair competition through dilution of the good will of their
brand, via the use of their BMW logo on a varnishing product from Verimark,
where Judge Stevens said:
\begin{quote}
\textit{``I fail to see how the use of the logo can affect the advertising value of the logo detrimentally. A mental association does not necessarily lead to bluring or tarnishing.''}
\end{quote}

Regardless however, Aroma Company may seek relief against Beaver Company, from a
High Court having jurisdiction, in the form of an interdict [section
34(3)(a)]\cite{rsa93_tm_act}, removal of the infringing mark or delivery of the
infringing goods [section 34(3)(b)]\cite{rsa93_tm_act}, damages or reasonable
royalties [sections 34(3)(c-d) and 34(4)]\cite{rsa93_tm_act}.

\section{Advise Aroma Company whether any formalities are necessary and/or advisable to protect its trade mark rights. [25]}
\label{sec:org359c643}

As per the provisions of [section 38(1)]\cite{rsa93_tm_act}, Clover Company may
use the TROPIC GLOW trade mark as a licensee with consent from the proprietor
Aroma Company. Permitted use of the TROPIC GLOW mark by Clover Company shall be
deemed as use by the proprietor for statutory purposes, [section
38(2)]\cite{rsa93_tm_act}. Moreover given that [section 38(3)]\cite{rsa93_tm_act}
pertains specifically to registration as a registered user in respect of ``all
or any of the goods or services'', implies that a license may be granted for use
on only some of the goods for the trade mark is registered, on condition that
such use is neither deceptive nor confusing. It would be advisable for Aroma
Company to explicitly articulate exactly which goods and / or services they
would like to appoint Clover Company as a licensee.\\

Aroma Company should be aware that they may be called upon by Clover Company to
institute infringement proceedings in respect of TROPIC GLOW, and should Aroma
Company not do so within two months after being so called upon, Clover Company
may independently institute such proceedings as if they were the proprietors of
the mark as per the provisions of [section 34]\cite{rsa93_tm_act}. In fulfilling
the licensee agreement, Aroma Company will be required to apply in writing to
the registrar in the prescribed manner, proposing the registration of Clover
Company as a registered user in respect of the TROPIC GLOW mark, furnishing the
registrar with the \begin{enumerate*}[label=(\roman*)]\item name and address of
Clover Company [section 38(6)(a)]\cite{rsa93_tm_act}, \item the existing and / or
proposed relationship between Aroma Company and Clover Company [section
38(6)(b)]\cite{rsa93_tm_act}, \item the goods or services in respect of which
Clover Company is to be registered as a registered user of TROPIC GLOW, [section
38(6)(c)]\cite{rsa93_tm_act} \end{enumerate*}.\\

Clover Company will be required to inform the public of the license agreement
through public notification and an ensure that it is recorded in the trade mark
register as a registered user, [section 38(3)-(7)]\cite{rsa93_tm_act}. This will
enable Clover Company to institute proceedings for infringement against use by
others of the registered trade mark TROPIC GLOW, in respect of the cosmetic
goods to which their license pertains, should Aroma Company fail to do so, as
per the provisions of [section 34]\cite{rsa93_tm_act}.\\

Should Aroma Company desire to have the license revoked or amended, they will
need to apply to the registrar in writing, in the prescribed manner [section
38(8)(a)]\cite{rsa93_tm_act}. Moreover the registrar shall cancel the license
should the mark be assigned as per [section 40]\cite{rsa93_tm_act}, unless the new
proprietor makes provisions for the extension of said license as per [section
38(6)]\cite{rsa93_tm_act}. Lastly, the license may be revoked, should the TROPIC
GLOW mark cease to be registered in respect of class 3 - cosmetics, [section
38(9)]\cite{rsa93_tm_act}.

\printbibliography
\end{document}